\clearpage\fancyhead[L]{\small KOMET.}\hypertarget{native-338-lower}{}\label{native-338-lower}
\begin{center}\textbf{KOMET.}\par\vspace{4pt}\hrule\vspace{4pt}\fontsize{10}{12}\selectfont Astronomisches Jahrbuch für 1814. Seite 254\,..\,256. Berlin 1811.\par\vspace{4pt}\hrule\vspace{4pt}\end{center}

\newcommand{\ObservedNumeratorMark}{\hbox{\vrule width0pt height2.291667bp depth0pt\special{pdf:bcontent}\special{pdf:code q 0 g 1.208333 0.833333 m 1.083333 0.833333 l 1.041667 1.000000 l 0.708333 1.000000 l 0.708333 0.875000 l 0.541667 0.833333 l 0.541667 0.708333 l 0.375000 0.666667 l 0.333333 0.500000 l 0.000000 0.500000 l 0.000000 0.000000 l 0.500000 0.000000 l 0.541667 0.166667 l 0.875000 0.166667 l 0.875000 0.291667 l 1.041667 0.333333 l 1.041667 0.500000 l 1.208333 0.541667 l h 0.000000 2.083333 m 0.000000 1.958333 l 0.125000 1.958333 l 0.166667 1.791667 l 1.041667 1.791667 l 1.083333 1.583333 l 1.375000 1.583333 l 1.375000 2.083333 l 1.208333 2.125000 l 1.208333 2.291667 l 0.166667 2.291667 l 0.166667 2.125000 l h f* Q}\special{pdf:econtent}\kern1.375000bp}}

\begin{flushright}Göttingen, den 10. September 1811.\end{flushright}
\par\indent Den Kometen haben wir hier zuerst am 22. August in der Abenddämmerung gesehen. Die Stadt schränkt jedoch die Sicht von unserer Sternwarte ein, sodass wir erst im September Beobachtungen anstellen konnten. Dazu habe ich die Abstände von Fixsternen verwendet, die sich mit dem Sextanten auf dem Stativ sehr gut beobachten lassen. Bisher konnten noch nicht alle Beobachtungen reduziert werden. Heute kann ich erst die folgenden drei mitteilen, die auf Abständen von $\alpha$ Lyrae und $\alpha$ Aurigae beruhen.
\Needspace{9\baselineskip}
\begin{center}\fontsize{10}{12}\selectfont\setlength{\tabcolsep}{4pt}\renewcommand{\arraystretch}{1.15}\begin{tabular}{r|c|c|c}
1811 & Mittlere Zeit & Rektaszension & Nördliche Deklination\\
\hline
Sept. 4 & $8^{\mathrm u}\,28^{\mathrm m}\,47^{\mathrm s}$ & $158^\circ\,25'\,24''$ & $39^\circ\,18'\,2''$\\
6 & $8\;48\;38$ & $160\;23\;16$ & $40\;14\;16$\\
7 & $8\;57\;6$ & $161\;26\;10$ & $40\;41\;54$\\
\end{tabular}\end{center}

\par\indent Anhand dieser und neuer Beobachtungen von Herrn von \textsc{Zach} habe ich meine vorläufigen parabolischen Elemente, die im 130. Stück unserer gelehrten Anzeigen [S. 327 d. B.] abgedruckt sind, wie folgt verbessert.
\Needspace{11\baselineskip}
\begin{center}\fontsize{10}{12}\selectfont\setlength{\tabcolsep}{4pt}\renewcommand{\arraystretch}{1.15}\begin{tabular}{l}
Durchgang durch das Perihel 1811 Sept. 12\\
\quad $5^{\mathrm u}\,21^{\mathrm m}\,15^{\mathrm s}$ mittlere Zeit in Göttingen\\
\end{tabular}\end{center}

\begin{center}\fontsize{10}{12}\selectfont\setlength{\tabcolsep}{4pt}\renewcommand{\arraystretch}{1.15}\begin{tabular}{llr}
Logarithmus des kleinsten Abstands & .\;.\;.\;.\;.\;.\;.\;.\;. & $0.017060$\\
Länge des Perihels & .\;.\;.\;.\;.\;.\;.\;.\;.\;.\;.\;.\;.\;.\;. & $75^\circ\,17'\,34''$\\
Länge des aufsteigenden Knotens & .\;.\;.\;.\;.\;.\;.\;.\;.\;. & $140\;24\;13$\\
Neigung der Bahn & .\;.\;.\;.\;.\;.\;.\;.\;.\;.\;.\;.\;.\;.\;. & $73\;7\;17$\\
\end{tabular}\end{center}
\noindent Bewegung rückläufig.

\Needspace{19\baselineskip}\SourceLeaf{339}\par\indent Daraus habe ich den geozentrischen Lauf für die nächsten 4 Monate wie unten angegeben berechnet.

\begin{center}\fontsize{10}{12}\selectfont\setlength{\tabcolsep}{4pt}\renewcommand{\arraystretch}{1.15}\begin{tabular}{rr|c|c|c|c}
\multicolumn{2}{c|}{Göttinger Zeit} & \shortstack{Rektaszension} & \shortstack{Nördliche\\Deklination} & \shortstack{Log. des Abstands\\von der Erde} & \shortstack{Helligkeit.}\\
\hline
1811 Sept. 12 & $5^{\mathrm u}$ & $166^\circ\,54'$ & $42^\circ\,58'$ & $0.2079$ & $0.355$\\
17 & $5$ & $173\;46$ & $45\;15$ & $0.1836$ & $0.394$\\
22 & $6$ & $182\;2$ & $47\;17$ & $0.1594$ & $0.432$\\
Okt. 2 & $7$ & $203\;24$ & $49\;34$ & $0.1160$ & $0.490$\\
12 & $8$ & $228\;45$ & $47\;15$ & $0.0911$ & $0.492$\\
22 & $9$ & $251\;23$ & $39\;56$ & $0.0963$ & $0.421$\\
Nov. 1 & $10$ & $267\;58$ & $30\;32$ & $0.1312$ & $0.310$\\
11 & $11$ & $279\;38$ & $22\;18$ & $0.1851$ & $0.208$\\
21 & $12$ & $288\;13$ & $14\;55$ & $0.2415$ & $0.138$\\
Dez. 1 & $13$ & $294\;54$ & $9\;42$ & $0.2976$ & $0.093$\\
11 & $14$ & $300\;24$ & $5\;49$ & $0.3487$ & $0.064$\\
21 & $15$ & $303\;22$ & $3\;0$ & $0.3848$ & $0.048$\\
\end{tabular}\end{center}

\par\indent Von der Mondfinsternis am 2. September war hier wegen sehr ungünstigen Wetters wenig zu sehen. Die Bedeckung von $\lambda$ \emph{Aquarii} während der Finsternis (im Jahrbuch etwa 1 Stunde zu früh angegeben) ließ sich aber sehr gut beobachten. Eintritt und Austritt erfolgten am verfinsterten, aber noch sichtbaren Mondrand. Der Austritt ließ sich umso besser beobachten, weil einer meiner Zuhörer die Bedeckung im Voraus berechnet hatte und wir deshalb darauf vorbereitet waren.

\par\indent Eintritt von $\lambda$ Aquarii, mittlere Zeit $10^{\mathrm u}\,16^{\mathrm m}\,17^{\mathrm s}4$. Austritt $10^{\mathrm u}\,51^{\mathrm m}\,40^{\mathrm s}5$.

\par\indent Auf Ihren Wunsch füge ich noch einen kurzen Zusatz zu meiner Theoria motus corporum coelestium bei.
\Needspace{16\baselineskip}\begin{center}\emph{Zusatz zu Artikel 90 und 100 der Theoria motus corporum coelestium, vom Verfasser.}\end{center}
\par\indent Bei der wichtigen Aufgabe, aus zwei Radiusvektoren und dem eingeschlossenen Winkel die elliptischen oder hyperbolischen Elemente zu bestimmen, habe ich mit großem Vorteil eine Hilfsgröße verwendet, $\xi$ für die Ellipse und $\zeta$ für die Hyperbel. Dafür habe ich dem Werk eine Tafel beigefügt. Die Tafel wurde nach einem dort angegebenen Kettenbruch berechnet, dessen vollständige Herleitung dort aber fehlt. Zu seiner theoretischen Entwicklung, die mit anderen Untersuchungen zusammenhängt, habe ich bisher noch keine Gelegenheit gefunden. Mancher wird daher hier gern einen anderen Weg finden, auf dem sich diese Hilfsgröße ebenso bequem hätte berechnen lassen.

\par\indent Wir haben (Art. 90)
\[
\xi=x-\frac{5}{6}+\frac{10}{9X}=\frac{xX-\dfrac{5}{6}X+\dfrac{10}{9}}{X}
\]
\par\indent Setzt man für $X$ die dort angegebene Reihe ein, so geht der Zähler dieses Bruchs leicht über in

\[
\begin{aligned}
&\frac{8}{105}xx\bigl(1+\frac{2\mathord{.}8}{9}x+\frac{3\mathord{.}8\mathord{.}10}{9\mathord{.}11}xx\\
&\qquad+\frac{4\mathord{.}8\mathord{.}10\mathord{.}12}{9\mathord{.}11\mathord{.}13}x^3+\frac{5\mathord{.}8\mathord{.}10\mathord{.}12\mathord{.}14}{9\mathord{.}11\mathord{.}13\mathord{.}15}x^4+\text{etc.}\bigr)
\end{aligned}
\]
\par\indent Setzt man also die Reihe
\[
1+\frac{2\mathord{.}8}{9}x+\frac{3\mathord{.}8\mathord{.}10}{9\mathord{.}11}xx+\text{etc.}=A
\]
\SourceLeaf{340}\noindent so ergibt sich
\[
xX-\frac{5}{6}X+\frac{10}{9}=\frac{8}{105}Axx
\]
\[
X=\frac{\dfrac{4}{3}\bigl(1-\dfrac{12}{175}Axx\bigr)}{1-\dfrac{6}{5}x}
\]
\[
\xi=\frac{\dfrac{\mathord{\raisebox{7bp}{\ObservedNumeratorMark}\,}2}{35}Axx\bigl(1-\dfrac{6}{5}x\bigr)}{1-\dfrac{12}{175}Axx}
\]
\noindent Nach dieser Formel kann man $\xi$ stets bequem und sicher berechnen. Für $\zeta$ braucht man nur $-z$ anstelle von $x$ einzusetzen.

\Needspace{8\baselineskip}\par\indent Ich bemerke nur noch, dass sich $A$ nach folgender Formel noch bequemer berechnen lässt

\[
\begin{aligned}
A={}&(1-x)^{-\frac{3}{2}}\bigl(1+\frac{1\mathord{.}5}{2\mathord{.}9}x+\frac{1\mathord{.}3\mathord{.}5\mathord{.}7}{2\mathord{.}4\mathord{.}9\mathord{.}11}xx\\
&\qquad+\frac{1\mathord{.}3\mathord{.}5\mathord{.}5\mathord{.}7\mathord{.}9}{2\mathord{.}4\mathord{.}6\mathord{.}9\mathord{.}11\mathord{.}13}x^3+\text{etc.}\bigr)
\end{aligned}
\]
\noindent Die Herleitung dieser Reihe aus der vorigen beruht jedoch auf Überlegungen, die hier nicht ausgeführt werden können.
